\section{Composition With Sibling LPs}
\label{sec:composition}

LP-203 is one of seven LPs in the ZAP-stack umbrella
(LP-200--LP-218). The composition is orthogonal: each LP solves
one concern, and they stack at well-defined seams.

\begin{table}[h]
\centering
\small
\begin{tabular}{@{}lp{9.5cm}@{}}
\toprule
\textbf{LP} & \textbf{Composition with LP-203} \\
\midrule
LP-022 & ZAP wire protocol. The buffer LP-203's GPU kernels
  read \cite{lp022}. \\
LP-077 & Round digest binding. The TupleHash256 transcript
  that feeds the GPU Merkle compute \cite{lp077}. \\
LP-132 & Quasar GPU execution adapter. The consensus-side hook
  this LP wires up \cite{lp132}. \\
LP-177 & Bridge PQ-only profile. Requires GPU on validators
  per the strict-PQ architectural decision \cite{lp177}. \\
LP-181 & Magnetar threshold SLH-DSA. CPU/GPU hybrid verify;
  this LP specs the crossover threshold \cite{lp181}. \\
LP-200 & ZAP Stack umbrella; the no-codec wire guarantee the
  GPU kernels rely on \cite{lp200}. \\
LP-201 & ZAP P2P / transport selector. DPDK and GPUDirect are
  ingest mechanisms underneath whichever transport
  \texttt{Pick()} chose for the peer \cite{lp201}. \\
LP-202 & ZAP pipelining and atomic unwind. The verify pipeline
  width LP-203's GPU kernels saturate \cite{lp202}. \\
LP-205 & Pipelined Quasar. The throughput model parameterized
  by GPU dispatch and mode \cite{lp205}. \\
LP-217 & Cert profile modes. Per-leg GPU dispatch lets the
  four modes share kernels \cite{lp217}. \\
LP-218 & ZAP P3Q PQ-rollup. P3Q precompile at EVM slot
  \texttt{0x012205}; Pulsar verify runs on the pure-Go/CPU path in
  production, with GPU dispatch on the path this LP specifies
  \emph{projected} (behind the composite-verify wrapper landing)
  \cite{lp218}. \\
\bottomrule
\end{tabular}
\caption{Composition of LP-203 with sibling LPs.}
\label{tab:composition}
\end{table}

\subsection{LP-201 transport selector: orthogonal to LP-203}

LP-201 defines a peer-to-transport selector \cite{lp201}. The
default is QUIC; LP-207 RDMA-IB is selectable for high-throughput
racks. LP-203's DPDK and GPUDirect contributions sit
\emph{underneath} the transport: they change how UDP packets
reach the process (kernel-bypass plus NIC DMA into UMA), not how
peers map to transports.

\paragraph{Composition contract.} DPDK and GPUDirect do not
register against \texttt{Transport.Pick()}. A GPUDirect-capable
NIC DMAs packets into UMA regardless of which transport
\texttt{Pick()} chose for the peer. The QUIC transport on a
DPDK + GPUDirect host gets its UDP packets via the
DPDK PMD; the LP-207 RDMA-IB transport on the same host gets
its frames via the GPUDirect MR. Both end up in the same UMA
buffer that the verify kernel reads.

\subsection{LP-205 pipelined Quasar: GPU dispatch is the
saturation source}

LP-205 \cite{lp205} parameterizes Quasar's verify pipeline
width by the underlying kernel throughput. With LP-203's
measured BLS verify at \SI{49197}~pairings/sec sustained on
a single GB10, a 4-validator host can run a 4-way pipeline
where each stage saturates a different GPU primitive:

\begin{itemize}[leftmargin=2em,nosep]
\item Stage 1: NIC ingest (DPDK or GPUDirect) -- bytes land
  in UMA
\item Stage 2: BLS pairing verify (GPU) -- the load-bearing
  measured kernel
\item Stage 3: Corona aggregate (GPU; once
  \texttt{lux-lattice.pc} ships) -- per-validator parallel
\item Stage 4: Merkle root reduction (GPU sha3-256 chain)
  -- $\log_2(N)$ kernel-launch chain
\end{itemize}

The stages share nothing except UMA pointers. Each stage's
output is a UMA address; the next stage reads from that
address. No \texttt{cudaMemcpy} between stages.

\subsection{LP-217 cert profile modes: GPU dispatch per leg}

LP-217 \cite{lp217} defines the operator-facing 4-mode
cert ladder: PQ-off (BLS), PQ-fast (BLS + Pulsar), PQ-strict
(BLS + Pulsar + Corona), PQ-heavy (BLS + Pulsar + Corona +
Magnetar). Each leg is verified by an independent GPU dispatch:

\begin{itemize}[leftmargin=2em,nosep]
\item BLS leg: \texttt{cuda\_op\_bls12\_381\_pairing\_eip2537}
  (\SI{19}{\micro\second} per pair sustained, measured)
\item Pulsar leg: the FIPS 204 verify wrapper (projected
  \SI{50}{\micro\second}--\SI{100}{\micro\second} per cert
  once the composite wrapper lands)
\item Corona leg: NTT batch on GPU + combine (Corona combine
  measured at \SI{0.082}{\micro\second} per share)
\item Magnetar leg: SLH-DSA 192f composite verify (projected
  $\sim$\SI{2}{\milli\second} at $N{=}64$ once the wrapper
  lands)
\end{itemize}

The four legs dispatch in parallel CUDA streams. A final CPU
step reads four boolean results from UMA and computes the
QuasarCert verdict. The PQ-strict round budget is set so the
slowest leg (Corona) determines the wall-clock; the other
three legs overlap.

\paragraph{Strict-PQ-* requires GPU.} The architectural
decision from LP-203 stands: strict-PQ-* modes drop the BLS
leg and the survivable-CPU path \emph{disappears}. Corona
aggregate over 64 validators takes
$\sim$\SI{350}{\milli\second} on CPU; the strict-PQ round
budget is tighter than that. The Lux Operator
(\texttt{luxfi/operator}) refuses to admit a validator to a
strict-PQ-enabled chain unless the validator's
\texttt{gpu.enabled} is true and the boot-time GPU probe
succeeds.

\subsection{LP-218 P3Q precompile at slot \texttt{0x012205}}

LP-218 \cite{lp218} specifies a single EVM
precompile that verifies a Pulsar threshold signature against a
group public key. In production the precompile runs on the
pure-Go/CPU verify path; there is no production GPU kernel for PQ
signature verify today. On a Blackwell-class validator, P3Q
\emph{would} dispatch to the Pulsar GPU verify wrapper once it lands
(projected; behind the composite-verify wrapper landing).

\paragraph{Gas cost contract.} P3Q gas cost is dimensioned for
worst-case CPU verify, not for GPU dispatch. Operators pay the
same gas whether the validator runs P3Q on CPU or GPU; the
incentive to deploy GPU validators is that the validator's
own wall-clock budget tightens. This decoupling preserves the
gas-cost spec's portability across hardware classes.

\subsection{Composition diagram}

\begin{verbatim}
+-----------------------------+    LP-200 (one byte stream)
|  ZAP buffer (UMA-backed)    |<-- LP-022 wire format
|  cudaMallocManaged          |    LP-077 round digest binds
|  shared across CPU + GPU    |
+-----------------------------+
            |
            | DPDK ingress  (LP-201 underneath, not on top of selector)
            | GPUDirect RDMA (same)
            v
+-----------------------------+
|  GPU verify kernels         |<-- LP-132 Quasar hook
|  (luxfi/accel, ABI v12)     |    LP-181 Magnetar SLH-DSA
|  - BLS12-381 EIP-2537       |    LP-205 pipelined Quasar
|  - keccak256 / sha3 / b3    |    LP-217 4-mode dispatch
|  - Pulsar / Corona NTT      |    LP-218 P3Q precompile
+-----------------------------+
            |
            | UMA fence
            v
+-----------------------------+
|  Consensus verdict          |
|  + Merkle root              |
|  + ZapDB commit             |
+-----------------------------+
\end{verbatim}

\section{Activation Marker}

\begin{verbatim}
activates:      2025-12-25T16:20:00-08:00
activates-unix: 1766708400
\end{verbatim}

GPU-native verify is the architecture from genesis of the new
final Lux network. No backwards compatibility with the pre-Quasar
Edition CPU-only verify path. The pre-Quasar network
(\num{2020}--\num{2025}) is a separate network out of scope.

\section{Backwards Compatibility}

None. The pre-Quasar Edition Lux network had no GPU-native
verify path. CPU verify remains as the fallback within
LP-203's profile (per \texttt{luxfi/accel} dispatcher) for
hosts without a GPU; in the strict-PQ-* modes, CPU fallback
is explicitly not survivable.

\section{Security Considerations}

\paragraph{No new cryptographic primitives.} LP-203 specifies
dispatch and zero-copy wire-up. The cryptographic primitives
verified (BLS12-381 EIP-2537, Ed25519, secp256k1, ML-DSA,
SLH-DSA, Corona) are unchanged from their FIPS or
EIP specifications \cite{fips203,fips204,fips205,eip2537}.

\paragraph{GPU-side fault model.} GPU kernels run with
warp-level data parallelism. A faulty kernel implementation
that returns ``verify OK'' on a bad signature would break the
verify guarantee. The dispatcher's per-kernel timeout plus
CPU fallback provides a fault-tolerant degrade path; the
kernel KAT vectors (FIPS test vectors per scheme) catch
implementation faults at boot time.

\paragraph{UMA security model.} \texttt{cudaMallocManaged}
memory is in the same address space as host memory; a kernel
that exceeds its bounds writes into host pages. The kernels
do not malloc; they only read from caller-supplied UMA
pointers. The host-side bounds check is on
\texttt{accel.BatchVerify*}'s parameter validation.

\paragraph{Side-channel surface.} GPU verify times are
data-dependent in principle (Miller loop, NTT, SHA). The
production hot path uses constant-time CPU implementations on
fallback. GPU implementations follow the same constant-time
discipline; KAT-vector tests cross-check.

\section{Future Work}

\begin{itemize}[leftmargin=2em,nosep]
\item \textbf{Composite signature-verify wrappers.} Close the
  16-of-76 stub gap in \texttt{plugin.cpp}. Once landed, every
  cell in \S\ref{sec:projections} becomes measurable.
\item \textbf{\texttt{lux-lattice.pc} install on Spark.}
  Unblocks lattice GPU dispatch.
\item \textbf{LP-204 GPU-shared resource scheduling.} One GPU
  per validator host serving multiple chain-VMs' verify
  queues fairly.
\item \textbf{LP-205 FHE GPU runtime.} Wires \texttt{luxfi/fhe}
  and \texttt{accel/ops\_fhe.go} into the consensus hot path
  for encrypted votes.
\item \textbf{LP-206 TEE attestation for GPU verify.}
  H100 / GB10 SEV-SNP attestation of the GPU verify kernels so
  GPU-computed Merkle roots and verdicts carry a remote-
  attestation receipt.
\item \textbf{CXL 3.0 racks.} Add cross-socket UMA semantics to
  the bench when CXL 3.0 hosts land in the bench cluster.
\end{itemize}
